% Generated draft: requires mathematical and source review.
\documentclass[11pt]{article}
\usepackage[margin=25mm]{geometry}
\usepackage{fontspec}
\setmainfont{Latin Modern Roman}
\usepackage{amsmath,amssymb,mathtools,unicode-math}
\setmathfont{Latin Modern Math}
\usepackage{enumitem,needspace,xurl,hyperref}
\setlength{\parindent}{0pt}
\setlength{\parskip}{5pt}
\setlength{\emergencystretch}{3em}
\newcommand{\N}{\mathbb{N}}
\newcommand{\Z}{\mathbb{Z}}
\newcommand{\Q}{\mathbb{Q}}
\newcommand{\R}{\mathbb{R}}
\newcommand{\C}{\mathbb{C}}
\newcommand{\F}{\mathbb{F}}
\newcommand{\A}{\mathbb{A}}
\newcommand{\PP}{\mathbb{P}}
\newcommand{\E}{\mathbb{E}}
\newcommand{\eps}{\varepsilon}
\newcommand{\dd}{\,\mathrm{d}}
\begin{document}
\section*{UnsolvedMath ID: 30000879}
\textbf{Source problem code:} catalogue entry
\subsection{Definitions and mathematical statement}
Let $G$ be a torsion-free nilpotent group of finite cohomological dimension $n$, and let $\alpha\in H^2(G,\C^\times)$. In the usual twisted-group-algebra convention, the coefficient action is trivial. Choose a normalized cocycle $a:G\times G\to\C^\times$ representing $\alpha$, so $a(g,h)a(gh,k)=a(h,k)a(g,hk)$ and $a(1,g)=a(g,1)=1$. The algebra $\C^\alpha G$ has $\C$-basis $\{u_g:g\in G\}$ and multiplication $u_gu_h=a(g,h)u_{gh}$; changing the representative cocycle gives an isomorphic algebra. Write $\operatorname{gl.dim}R$ for global dimension, conventionally the supremum of projective dimensions of left $R$-modules, and interpret cohomological dimension conventionally as $\operatorname{cd}_{\Z}G=\operatorname{pd}_{\Z G}\Z$ unless the report specifies another coefficient ring. Suppose \[1\leq r\leq n,\qquad \operatorname{gl.dim}(\C^\alpha G)=r.\] Must there exist a subgroup $H\leq G$ such that \[\operatorname{res}^G_H\alpha=0\in H^2(H,\C^\times),\qquad\operatorname{gl.dim}(\C H)=r?\] Trivial restriction is a cohomological assertion: it means that $a|_{H\times H}$ is a coboundary, for example $a(h,k)=b(h)b(k)/b(hk)$ for some $b:H\to\C^\times$, not necessarily that the originally chosen cocycle is pointwise one. No finite generation of $G$, normality or finite index of $H$, or equality of $H$ with a maximal such subgroup is required by the extract.\par The supplied source fields conflict: the original and clean/current statements give this complete question, and the corrected\_verified note says its cohomology/twisted-algebra/global-dimension notation was restored. The 2026-08-21 literature assessment instead describes an extraction containing only the word conjecture with no recoverable claim and says this was retained as open for that reason. Both descriptions are preserved; the assessment is not treated as a verification of the displayed conjecture's current mathematical status.
\subsection{Short English statement}
If a twisted complex group algebra of a torsion-free nilpotent group has global dimension $r$, ask whether the twist becomes cohomologically trivial on some subgroup whose ordinary group algebra has exactly that global dimension. The literature summary conflicts with the supplied complete statement.
\subsection{Sources}
\begin{itemize}\item Local UnsolvedMath registry, permanent ID 30000879, source problem catalogue entry, \emph{Trivial Restrictions of Twists in Nilpotent Group Algebras}, proposed year 2007; no individual proposer supplied.\item Supplied citation: \emph{Mini-Workshop: Arithmetik von Gruppenringen} (2008), Oberwolfach Reports, DOI 10.4171/owr/2007/55, pp.\ 3209--3240: \url{https://doi.org/10.4171/owr/2007/55}. Supplied literature label: original Oberwolfach report.\end{itemize}
\end{document}
